% GENERATED FILE. Edit canonical lessons, then run npm run generate:books.
% canonical-source-hash: fnv1a64:330132d5098c9e2b
% canonical-lessons: BN-C441-atongko, BN-C441-atke-otha, BN-C441-shiure-otha, BN-C441-ghabrano, BN-C441-dhorphor-kora

\chapter{A Sudden Fright}
\label{ch:bn-a-sudden-fright}
\hlchaptermodality{441}

\begin{chapteropening}
\textbf{By the end of this chapter:} \emph{I can say that something frightened me.}

Complete the last lesson of chapter 441: I can say that something frightened me.
\end{chapteropening}

\section[\texorpdfstring{\bn{আতঙ্ক} (ātôngkô)}{আতঙ্ক (ātôngkô)}]{\bn{আতঙ্ক} (ātôngkô) — terror; panic}

\label{lesson:BN-C441-atongko}



\begin{quote}
Before the new one: say the Bengali for to be surprised, then the Bengali for to be startled.
\end{quote}

\subsection*{You'll want to know: \bn{আতঙ্ক}}
\textbf{\bn{আতঙ্ক}} — \emph{ātôngkô} — \textquotedblleft{}terror; panic\textquotedblright{}.

\begin{quote}
\bn{আগুন} \bn{দেখে} \bn{গ্রামে} \bn{আতঙ্ক} \bn{ছড়াল।}
\end{quote}

\begin{quote}
\emph{āgun dekhe grāme ātôngkô chhôṛālô.}
\end{quote}

\begin{quote}
Seeing the fire, terror spread through the village.
\end{quote}

\bn{আতঙ্কে}, in terror: \bn{আতঙ্কে} \bn{কাঁপা}, to shake with terror.

\begin{cousinweb}
\textbf{\bn{আতঙ্ক}} (\emph{ātôngkô}) is a tatsama, Sanskrit \emph{ātaṅka}, which Monier-Williams gives for sickness and pain and also for fear; Bengali keeps the fear.
\end{cousinweb}

\begin{grammarlens}[title={What you've built}]
One more word for describing what you lived and felt.
\end{grammarlens}

\subsection*{Your turn}
\begin{itemize}
\raggedright
  \item \emph{Say it:} \emph{ātôngkô}
  \item \emph{Say it:} \emph{ātôngkô}, once more
  \item \emph{Say it:} \emph{āgun dekhe grāme ātôngkô chhôṛālô.} — Seeing the fire, terror spread through the village.
  \item \emph{Recall:} say the Bengali for to be surprised, then the Bengali for to be startled, then say \emph{ātôngkô} again
\end{itemize}

\begin{tcolorbox}[breakable,colback=teal!4,colframe=teal!35!black,title={Before you move on}]
What does \bn{আতঙ্ক} mean? (\textquotedblleft{}Terror; panic\textquotedblright{}.) Say it once more.
\end{tcolorbox}
\section[\texorpdfstring{\bn{আঁতকে} \bn{ওঠা} (ā̃tke oṭhā)}{আঁতকে ওঠা (ā̃tke oṭhā)}]{\bn{আঁতকে} \bn{ওঠা} (ā̃tke oṭhā) — to recoil in fright; to gasp with alarm}

\label{lesson:BN-C441-atke-otha}



\begin{quote}
Before the new one: say the Bengali for to be startled, then the Bengali for terror.
\end{quote}

\subsection*{You'll want to know: \bn{আঁতকে} \bn{ওঠা}}
\textbf{\bn{আঁতকে} \bn{ওঠা}} — \emph{ā̃tke oṭhā} — \textquotedblleft{}to recoil in fright; to gasp with alarm\textquotedblright{}.

\begin{quote}
\bn{সাপটা} \bn{দেখে} \bn{মীরা} \bn{আঁতকে} \bn{উঠল।}
\end{quote}

\begin{quote}
\emph{sāpṭā dekhe mīrā ā̃tke uṭhlô.}
\end{quote}

\begin{quote}
Mira recoiled in fright at the sight of the snake.
\end{quote}

Past (he/she …): \bn{তিনি} \bn{খবর} \bn{শুনে} \bn{আঁতকে} \bn{উঠলেন}, \emph{tini khôbôr shune ā̃tke uṭhlen}.

\begin{cousinweb}
\textbf{\bn{আঁতকে} \bn{ওঠা}} (\emph{ā̃tke oṭhā}) is \emph{ā̃tke}, the -\bn{এ} form of a verb of sudden fright whose origin is not settled, with \emph{oṭhā}, \textquotedblleft{}to rise\textquotedblright{} (CDIAL 1793).
\end{cousinweb}

\begin{grammarlens}[title={What you've built}]
One more word for describing what you lived and felt.
\end{grammarlens}

\subsection*{Your turn}
\begin{itemize}
\raggedright
  \item \emph{Say it:} \emph{ā̃tke oṭhā}
  \item \emph{Say it:} \emph{ā̃tke oṭhā}, once more
  \item \emph{Say it:} \emph{sāpṭā dekhe mīrā ā̃tke uṭhlô.} — Mira recoiled in fright at the sight of the snake.
  \item \emph{Recall:} say the Bengali for to be startled, then the Bengali for terror, then say \emph{ā̃tke oṭhā} again
\end{itemize}

\begin{tcolorbox}[breakable,colback=teal!4,colframe=teal!35!black,title={Before you move on}]
What does \bn{আঁতকে} \bn{ওঠা} mean? (\textquotedblleft{}To recoil in fright; to gasp with alarm\textquotedblright{}.) Say it once more.
\end{tcolorbox}
\section[\texorpdfstring{\bn{শিউরে} \bn{ওঠা} (shiure oṭhā)}{শিউরে ওঠা (shiure oṭhā)}]{\bn{শিউরে} \bn{ওঠা} (shiure oṭhā) — to shudder}

\label{lesson:BN-C441-shiure-otha}



\begin{quote}
Before the new one: say the Bengali for terror, then the Bengali for to recoil in fright.
\end{quote}

\subsection*{You'll want to know: \bn{শিউরে} \bn{ওঠা}}
\textbf{\bn{শিউরে} \bn{ওঠা}} — \emph{shiure oṭhā} — \textquotedblleft{}to shudder\textquotedblright{}.

\begin{quote}
\bn{গল্পটা} \bn{শুনে} \bn{আমি} \bn{শিউরে} \bn{উঠলাম।}
\end{quote}

\begin{quote}
\emph{gôlpôṭā shune āmi shiure uṭhlām.}
\end{quote}

\begin{quote}
I shuddered when I heard the story.
\end{quote}

Past (he/she …): \bn{তিনি} \bn{ছবিটা} \bn{দেখে} \bn{শিউরে} \bn{উঠলেন}, \emph{tini chhôbiṭā dekhe shiure uṭhlen}.

\begin{cousinweb}
\textbf{\bn{শিউরে} \bn{ওঠা}} (\emph{shiure oṭhā}) is \emph{shiure}, the -\bn{এ} form of a verb that is probably a softened form of the Bengali \emph{siharā}, of hair standing on end, which CDIAL files with \emph{siharāna}, \textquotedblleft{}to startle\textquotedblright{}, under Sanskrit \emph{śikhara}, \textquotedblleft{}a peak\textquotedblright{}, a word the lexicographers also use for bristling hair (CDIAL 12435), with \emph{oṭhā}, \textquotedblleft{}to rise\textquotedblright{} (CDIAL 1793).
\end{cousinweb}

\begin{grammarlens}[title={What you've built}]
One more word for describing what you lived and felt.
\end{grammarlens}

\subsection*{Your turn}
\begin{itemize}
\raggedright
  \item \emph{Say it:} \emph{shiure oṭhā}
  \item \emph{Say it:} \emph{shiure oṭhā}, once more
  \item \emph{Say it:} \emph{gôlpôṭā shune āmi shiure uṭhlām.} — I shuddered when I heard the story.
  \item \emph{Recall:} say the Bengali for terror, then the Bengali for to recoil in fright, then say \emph{shiure oṭhā} again
\end{itemize}

\begin{tcolorbox}[breakable,colback=teal!4,colframe=teal!35!black,title={Before you move on}]
What does \bn{শিউরে} \bn{ওঠা} mean? (\textquotedblleft{}To shudder\textquotedblright{}.) Say it once more.
\end{tcolorbox}
\section[\texorpdfstring{\bn{ঘাবড়ানো} (ghābṛāno)}{ঘাবড়ানো (ghābṛāno)}]{\bn{ঘাবড়ানো} (ghābṛāno) — to get flustered; to lose one's nerve}

\label{lesson:BN-C441-ghabrano}



\begin{quote}
Before the new one: say the Bengali for to recoil in fright, then the Bengali for to shudder.
\end{quote}

\subsection*{You'll want to know: \bn{ঘাবড়ানো}}
\textbf{\bn{ঘাবড়ানো}} — \emph{ghābṛāno} — \textquotedblleft{}to get flustered; to lose one's nerve\textquotedblright{}.

\begin{quote}
\bn{প্রশ্নটা} \bn{শুনে} \bn{অরুণ} \bn{ঘাবড়ে} \bn{গেল।}
\end{quote}

\begin{quote}
\emph{prôshnôṭā shune ôrun ghābṛe gelô.}
\end{quote}

\begin{quote}
Arun got flustered when he heard the question.
\end{quote}

Past (he/she …): \bn{তিনি} \bn{ঘাবড়ালেন}, \emph{tini ghābṛālen}.

Usually in its -\bn{এ} form with \bn{যাওয়া}: \bn{ঘাবড়ে} \bn{যাওয়া}.

\begin{cousinweb}
\textbf{\bn{ঘাবড়ানো}} (\emph{ghābṛāno}) is listed by CDIAL for Bengali, as being confused and flustered, under a reconstructed \emph{ghabbhara}, \textquotedblleft{}confused\textquotedblright{} (CDIAL 4427), the entry where CDIAL also files the Hindi word for being startled.
\end{cousinweb}

\begin{grammarlens}[title={What you've built}]
One more word for describing what you lived and felt.
\end{grammarlens}

\subsection*{Your turn}
\begin{itemize}
\raggedright
  \item \emph{Say it:} \emph{ghābṛāno}
  \item \emph{Say it:} \emph{ghābṛāno}, once more
  \item \emph{Say it:} \emph{prôshnôṭā shune ôrun ghābṛe gelô.} — Arun got flustered when he heard the question.
  \item \emph{Recall:} say the Bengali for to recoil in fright, then the Bengali for to shudder, then say \emph{ghābṛāno} again
\end{itemize}

\begin{tcolorbox}[breakable,colback=teal!4,colframe=teal!35!black,title={Before you move on}]
What does \bn{ঘাবড়ানো} mean? (\textquotedblleft{}To get flustered; to lose one's nerve\textquotedblright{}.) Say it once more.
\end{tcolorbox}
\section[\texorpdfstring{\bn{ধড়ফড়} \bn{করা} (dhôṛphôṛ kôrā)}{ধড়ফড় করা (dhôṛphôṛ kôrā)}]{\bn{ধড়ফড়} \bn{করা} (dhôṛphôṛ kôrā) — to pound; to flutter (of the heart)}

\label{lesson:BN-C441-dhorphor-kora}



\begin{quote}
Before the new one: say the Bengali for to shudder, then the Bengali for to get flustered.
\end{quote}

\subsection*{You'll want to know: \bn{ধড়ফড়} \bn{করা}}
\textbf{\bn{ধড়ফড়} \bn{করা}} — \emph{dhôṛphôṛ kôrā} — \textquotedblleft{}to pound; to flutter (of the heart)\textquotedblright{}.

\begin{quote}
\bn{ভয়ে} \bn{আমার} \bn{বুক} \bn{ধড়ফড়} \bn{করছিল।}
\end{quote}

\begin{quote}
\emph{bhôye āmār buk dhôṛphôṛ kôrchhilô.}
\end{quote}

\begin{quote}
My heart was pounding with fear.
\end{quote}

Past (he/she …): \bn{সেই} \bn{রোগীর} \bn{বুক} \bn{ধড়ফড়} \bn{করল}, \emph{sei rogīr buk dhôṛphôṛ kôrlô}.

The heart is \bn{বুক} (taught), the chest: \bn{বুক} \bn{ধড়ফড়} \bn{করা}.

\begin{cousinweb}
\textbf{\bn{ধড়ফড়} \bn{করা}} (\emph{dhôṛphôṛ kôrā}) is \emph{dhôṛphôṛ}, an imitative word for a quick, fluttering beat, with \emph{kôrā}, \textquotedblleft{}to do\textquotedblright{} (CDIAL 2814).
\end{cousinweb}

\begin{grammarlens}[title={What you've built}]
One more word for describing what you lived and felt.
\end{grammarlens}

\subsection*{Your turn}
\begin{itemize}
\raggedright
  \item \emph{Say it:} \emph{dhôṛphôṛ kôrā}
  \item \emph{Say it:} \emph{dhôṛphôṛ kôrā}, once more
  \item \emph{Say it:} \emph{bhôye āmār buk dhôṛphôṛ kôrchhilô.} — My heart was pounding with fear.
  \item \emph{Recall:} say the Bengali for to shudder, then the Bengali for to get flustered, then say \emph{dhôṛphôṛ kôrā} again
\end{itemize}

\begin{tcolorbox}[breakable,colback=teal!4,colframe=teal!35!black,title={Before you move on}]
What does \bn{ধড়ফড়} \bn{করা} mean? (\textquotedblleft{}To pound; to flutter (of the heart)\textquotedblright{}.) Say it once more.
\end{tcolorbox}
